\documentclass[10pt,a4paper]{amsart}
\usepackage{amsmath,amssymb,amsthm,graphicx}
\usepackage{verbatim}
\newenvironment{reference}{\comment}{\endcomment}
%\newenvironment{reference}{}{}
\newenvironment{slogan}{\comment}{\endcomment}
\newenvironment{history}{\comment}{\endcomment}
% For commutative diagrams we use Xy-pic
\usepackage[all]{xy}

% We use 2cell for 2-commutative diagrams.
\xyoption{2cell}
% Theorem environments.
%
\theoremstyle{plain}
\newtheorem{theorem}[subsection]{Theorem}
\newtheorem{proposition}[subsection]{Proposition}
\newtheorem{lemma}[subsection]{Lemma}

\theoremstyle{definition}
\newtheorem{definition}[subsection]{Definition}
\newtheorem{example}[subsection]{Example}
\newtheorem{exercise}[subsection]{Exercise}
\newtheorem{situation}[subsection]{Situation}

\theoremstyle{remark}
\newtheorem{remark}[subsection]{Remark}
\newtheorem{remarks}[subsection]{Remarks}

\numberwithin{equation}{subsection}

% Macros
%
\def\lim{\mathop{\mathrm{lim}}\nolimits}
\def\colim{\mathop{\mathrm{colim}}\nolimits}
\def\Spec{\mathop{\mathrm{Spec}}}
\def\Hom{\mathop{\mathrm{Hom}}\nolimits}
\def\Ext{\mathop{\mathrm{Ext}}\nolimits}
\def\SheafHom{\mathop{\mathcal{H}\!\mathit{om}}\nolimits}
\def\SheafExt{\mathop{\mathcal{E}\!\mathit{xt}}\nolimits}
\def\Sch{\mathit{Sch}}
\def\Mor{\mathop{\mathrm{Mor}}\nolimits}
\def\Ob{\mathop{\mathrm{Ob}}\nolimits}
\def\Sh{\mathop{\mathit{Sh}}\nolimits}
\def\NL{\mathop{N\!L}\nolimits}
\def\CH{\mathop{\mathrm{CH}}\nolimits}
\def\proetale{{pro\text{-}\acute{e}tale}}
\def\etale{{\acute{e}tale}}
\def\QCoh{\mathit{QCoh}}
\def\Ker{\mathop{\mathrm{Ker}}}
\def\Im{\mathop{\mathrm{Im}}}
\def\Coker{\mathop{\mathrm{Coker}}}
\def\Coim{\mathop{\mathrm{Coim}}}

% Boxtimes
%
\DeclareMathSymbol{\boxtimes}{\mathbin}{AMSa}{"02}

%
% Macros for moduli stacks/spaces
%
\def\QCohstack{\mathcal{QC}\!\mathit{oh}}
\def\Cohstack{\mathcal{C}\!\mathit{oh}}
\def\Spacesstack{\mathcal{S}\!\mathit{paces}}
\def\Quotfunctor{\mathrm{Quot}}
\def\Hilbfunctor{\mathrm{Hilb}}
\def\Curvesstack{\mathcal{C}\!\mathit{urves}}
\def\Polarizedstack{\mathcal{P}\!\mathit{olarized}}
\def\Complexesstack{\mathcal{C}\!\mathit{omplexes}}
% \Pic is the operator that assigns to X its picard group, usage \Pic(X)
% \Picardstack_{X/B} denotes the Picard stack of X over B
% \Picardfunctor_{X/B} denotes the Picard functor of X over B
\def\Pic{\mathop{\mathrm{Pic}}\nolimits}
\def\Picardstack{\mathcal{P}\!\mathit{ic}}
\def\Picardfunctor{\mathrm{Pic}}
\def\Deformationcategory{\mathcal{D}\!\mathit{ef}}
\usepackage{cleveref}
\crefname{originallemma}{original lemma}{original lemmas}
\Crefname{originallemma}{Original lemma}{Original lemmas}
\makeatletter
\newlabel{algebra-lemma-characterize-separable-algebraic-field-extensions}{{090W}{source}{}{sourcepointer-1}{}}
\newlabel{algebra-lemma-characterize-separable-algebraic-field-extensions@cref}{{[originallemma][0][]090W}{[1][0][]source}{}{}{}}
\makeatother
\crefname{originallemma}{original lemma}{original lemmas}
\Crefname{originallemma}{Original lemma}{Original lemmas}
\makeatletter
\newlabel{coherent-lemma-flat-base-change-cohomology}{{02KH}{source}{}{sourcepointer-2}{}}
\newlabel{coherent-lemma-flat-base-change-cohomology@cref}{{[originallemma][0][]02KH}{[1][0][]source}{}{}{}}
\makeatother
\crefname{originallemma}{original lemma}{original lemmas}
\Crefname{originallemma}{Original lemma}{Original lemmas}
\makeatletter
\newlabel{coherent-lemma-proper-over-affine-cohomology-finite}{{02O6}{source}{}{sourcepointer-3}{}}
\newlabel{coherent-lemma-proper-over-affine-cohomology-finite@cref}{{[originallemma][0][]02O6}{[1][0][]source}{}{}{}}
\makeatother
\crefname{originallemma}{original lemma}{original lemmas}
\Crefname{originallemma}{Original lemma}{Original lemmas}
\makeatletter
\newlabel{coherent-lemma-relative-affine-cohomology}{{089W}{source}{}{sourcepointer-4}{}}
\newlabel{coherent-lemma-relative-affine-cohomology@cref}{{[originallemma][0][]089W}{[1][0][]source}{}{}{}}
\makeatother
\crefname{originalproposition}{original proposition}{original propositions}
\Crefname{originalproposition}{Original proposition}{Original propositions}
\makeatletter
\newlabel{cohomology-proposition-vanishing-Noetherian}{{02UZ}{source}{}{sourcepointer-5}{}}
\newlabel{cohomology-proposition-vanishing-Noetherian@cref}{{[originalproposition][0][]02UZ}{[1][0][]source}{}{}{}}
\makeatother
\crefname{originalremark}{original remark}{original remarks}
\Crefname{originalremark}{Original remark}{Original remarks}
\makeatletter
\newlabel{duality-remark-coherent-duality-proper-over-field}{{0FVX}{source}{}{sourcepointer-6}{}}
\newlabel{duality-remark-coherent-duality-proper-over-field@cref}{{[originalremark][0][]0FVX}{[1][0][]source}{}{}{}}
\makeatother
\crefname{originallemma}{original lemma}{original lemmas}
\Crefname{originallemma}{Original lemma}{Original lemmas}
\makeatletter
\newlabel{lemma-Garel-map-frobenius-smooth-char-p}{{0FW2}{source}{}{sourcepointer-7}{}}
\newlabel{lemma-Garel-map-frobenius-smooth-char-p@cref}{{[originallemma][0][]0FW2}{[1][0][]source}{}{}{}}
\makeatother
\crefname{originallemma}{original lemma}{original lemmas}
\Crefname{originallemma}{Original lemma}{Original lemmas}
\makeatletter
\newlabel{lemma-bottom-part-degenerates}{{0FW5}{source}{}{sourcepointer-8}{}}
\newlabel{lemma-bottom-part-degenerates@cref}{{[originallemma][0][]0FW5}{[1][0][]source}{}{}{}}
\makeatother
\crefname{originallemma}{original lemma}{original lemmas}
\Crefname{originallemma}{Original lemma}{Original lemmas}
\makeatletter
\newlabel{lemma-duality-hodge}{{0FW4}{source}{}{sourcepointer-9}{}}
\newlabel{lemma-duality-hodge@cref}{{[originallemma][0][]0FW4}{[1][0][]source}{}{}{}}
\makeatother
\crefname{originallemma}{original lemma}{original lemmas}
\Crefname{originallemma}{Original lemma}{Original lemmas}
\makeatletter
\newlabel{limits-lemma-descend-dimension-d}{{0EY2}{source}{}{sourcepointer-10}{}}
\newlabel{limits-lemma-descend-dimension-d@cref}{{[originallemma][0][]0EY2}{[1][0][]source}{}{}{}}
\makeatother
\crefname{originallemma}{original lemma}{original lemmas}
\Crefname{originallemma}{Original lemma}{Original lemmas}
\makeatletter
\newlabel{limits-lemma-descend-finite-presentation}{{01ZM}{source}{}{sourcepointer-11}{}}
\newlabel{limits-lemma-descend-finite-presentation@cref}{{[originallemma][0][]01ZM}{[1][0][]source}{}{}{}}
\makeatother
\crefname{originallemma}{original lemma}{original lemmas}
\Crefname{originallemma}{Original lemma}{Original lemmas}
\makeatletter
\newlabel{limits-lemma-descend-smooth}{{0C0C}{source}{}{sourcepointer-12}{}}
\newlabel{limits-lemma-descend-smooth@cref}{{[originallemma][0][]0C0C}{[1][0][]source}{}{}{}}
\makeatother
\crefname{originallemma}{original lemma}{original lemmas}
\Crefname{originallemma}{Original lemma}{Original lemmas}
\makeatletter
\newlabel{limits-lemma-eventually-proper}{{081F}{source}{}{sourcepointer-13}{}}
\newlabel{limits-lemma-eventually-proper@cref}{{[originallemma][0][]081F}{[1][0][]source}{}{}{}}
\makeatother
\crefname{originallemma}{original lemma}{original lemmas}
\Crefname{originallemma}{Original lemma}{Original lemmas}
\makeatletter
\newlabel{limits-lemma-higher-direct-images-zero-above-dimension-fibre}{{0E7D}{source}{}{sourcepointer-14}{}}
\newlabel{limits-lemma-higher-direct-images-zero-above-dimension-fibre@cref}{{[originallemma][0][]0E7D}{[1][0][]source}{}{}{}}
\makeatother
\crefname{originallemma}{original lemma}{original lemmas}
\Crefname{originallemma}{Original lemma}{Original lemmas}
\makeatletter
\newlabel{morphisms-lemma-triangle-differentials}{{01UX}{source}{}{sourcepointer-15}{}}
\newlabel{morphisms-lemma-triangle-differentials@cref}{{[originallemma][0][]01UX}{[1][0][]source}{}{}{}}
\makeatother
\crefname{originallemma}{original lemma}{original lemmas}
\Crefname{originallemma}{Original lemma}{Original lemmas}
\makeatletter
\newlabel{varieties-lemma-proper-geometrically-reduced-global-sections}{{0BUG}{source}{}{sourcepointer-16}{}}
\newlabel{varieties-lemma-proper-geometrically-reduced-global-sections@cref}{{[originallemma][0][]0BUG}{[1][0][]source}{}{}{}}
\makeatother
\crefname{originallemma}{original lemma}{original lemmas}
\Crefname{originallemma}{Original lemma}{Original lemmas}
\makeatletter
\newlabel{varieties-lemma-smooth-geometrically-normal}{{056T}{source}{}{sourcepointer-17}{}}
\newlabel{varieties-lemma-smooth-geometrically-normal@cref}{{[originallemma][0][]056T}{[1][0][]source}{}{}{}}
\makeatother
\title{The top differential of the Hodge-to-de Rham page vanishes for a smooth proper equidimensional scheme over a field}
\author{The Stacks Project Authors}
\date{stacks-project:derham.tex:local-snapshot-20261012; GNU FDL 1.2 or later}
\begin{document}
\maketitle
\noindent\textit{Editorial note.} All delivered mathematical text is the author's own, copied from the pinned Stacks Project chapter file derham.tex; no statement or proof step was rewritten, reordered or removed.\par\medskip
\noindent\textit{Editorial note.} Preamble provenance: the pinned chapter file opens with \textbackslash{}input\{preamble\} and carries no preamble text of its own, so the project's single Stacks-wide preamble staged at sources/stacks-derham/preamble.tex is excerpted verbatim as three spans, lines 8-12 (verbatim package and the project's own reference/slogan/history display environments), lines 14-18 (Xy-pic setup) and lines 154-217 (theorem declarations and mathematical macros). The wrapper class and the amsmath/amssymb/amsthm/graphicx base packages are supplied by the assembler. No macro was invented or redefined.\par\medskip
\noindent\textit{Editorial note.} Same-chapter pointers are located at their \textbackslash{}label lines inside the pinned chapter file; cross-chapter pointers are located in the project's own tag list staged in the snapshot (sources/stacks-chow/tags.txt), at the verbatim lines carrying their canonical tags.\par\medskip
\noindent\textit{Allowed standard prerequisites (editorial).}\begin{itemize}
\item The bottom part degeneration and the coherent duality remark, and why they do not suffice. Same chapter derham.tex, label lemma-bottom-part-degenerates (cited at proof line 4588) and Duality for Schemes, remark duality-remark-coherent-duality-proper-over-field (cited at proof line 4592). The proof opens by explaining that the tempting duality argument fails because the maps \$\textbackslash{}mathcal\{O\}\_X \textbackslash{}to \textbackslash{}Omega\textasciicircum{}1\_\{X/k\}\$ and \$\textbackslash{}Omega\textasciicircum{}\{d-1\}\_\{X/k\} \textbackslash{}to \textbackslash{}Omega\textasciicircum{}d\_\{X/k\}\$ are not \$\textbackslash{}mathcal\{O\}\_X\$-linear, so functoriality cannot identify the two maps as duals.
\item Smoothness, properness and the identification of the base field with the global functions. Varieties, labels varieties-lemma-smooth-geometrically-normal (invoked at proof line 4600) and varieties-lemma-proper-geometrically-reduced-global-sections (invoked at proof line 4601), Algebra, label algebra-lemma-characterize-separable-algebraic-field-extensions (invoked at proof line 4604) and Morphisms, label morphisms-lemma-triangle-differentials (invoked at proof line 4607). They reduce to an irreducible \$X\$ with \$H\textasciicircum{}0(X, \textbackslash{}mathcal\{O\}\_X) = k\$, replacing \$k\$ by a finite separable extension and using \$\textbackslash{}Omega\_\{k'/k\} = 0\$.
\item Hodge duality in the top degree. Same chapter derham.tex, label lemma-duality-hodge (invoked at proof lines 4612 and 4642). It gives \$\textbackslash{}dim H\textasciicircum{}d(X, \textbackslash{}Omega\textasciicircum{}d\_\{X/k\}) = 1\$ and shows that the corresponding group for \$X\textasciicircum{}\{(p)\}\$ is nonzero, which is what makes the surjectivity argument for \$\textbackslash{}Theta\textasciicircum{}d\$ sufficient.
\item The relative Frobenius in characteristic \$p\$ and the Garel maps attached to it. Same chapter derham.tex, label lemma-Garel-map-frobenius-smooth-char-p (invoked at proof line 4616). It provides the properties of \$F\_\{X/k\} : X \textbackslash{}to X\textasciicircum{}\{(p)\}\$ and of the maps \$\textbackslash{}Theta\textasciicircum{}\{d-1\}, \textbackslash{}Theta\textasciicircum{}d\$ used in the commutative diagram; in particular \$\textbackslash{}Theta\textasciicircum{}\{d-1\} = 0\$ and \$\textbackslash{}Theta\textasciicircum{}d\$ is surjective.
\item Finiteness of the Frobenius, right exactness in top degree, and the spreading out of a smooth proper family. Cohomology of Schemes, labels coherent-lemma-relative-affine-cohomology (invoked at proof line 4629) and coherent-lemma-proper-over-affine-cohomology-finite (invoked at proof line 4662); Cohomology, label cohomology-proposition-vanishing-Noetherian (invoked at proof line 4640); Limits, labels limits-lemma-descend-finite-presentation, limits-lemma-descend-smooth, limits-lemma-descend-dimension-d and limits-lemma-eventually-proper (invoked at proof lines 4658-4660) and limits-lemma-higher-direct-images-zero-above-dimension-fibre (invoked at proof line 4681); Cohomology of Schemes, label coherent-lemma-flat-base-change-cohomology (invoked at proof line 4667). They make the horizontal arrows isomorphisms (finiteness of \$F\$), make \$H\textasciicircum{}d\$ right exact (vanishing above \$d\$), spread \$X\$ out to a smooth proper family over a finite type \$\textbackslash{}mathbf\{Z\}\$-algebra, make its cohomology finite free after localization and identify the special fibres by flat base change.
\end{itemize}\medskip
\ifdefined\hypertarget\hypertarget{sourcepointer-1}{}\fi
\noindent\textit{Source reference.} Original lemma 090W: sources/stacks-chow/tags.txt, source lines 11424--11424. Navigation pointer for \textbackslash{}ref\{algebra-lemma-characterize-separable-algebraic-field-extensions\} used inside the counted span (line 4605); the argument that consumes it is retained in the span.\par\medskip
\ifdefined\hypertarget\hypertarget{sourcepointer-2}{}\fi
\noindent\textit{Source reference.} Original lemma 02KH: sources/stacks-chow/tags.txt, source lines 3372--3372. Navigation pointer for \textbackslash{}ref\{coherent-lemma-flat-base-change-cohomology\} used inside the counted span (line 4668); the argument that consumes it is retained in the span.\par\medskip
\ifdefined\hypertarget\hypertarget{sourcepointer-3}{}\fi
\noindent\textit{Source reference.} Original lemma 02O6: sources/stacks-chow/tags.txt, source lines 3505--3505. Navigation pointer for \textbackslash{}ref\{coherent-lemma-proper-over-affine-cohomology-finite\} used inside the counted span (line 4663); the argument that consumes it is retained in the span.\par\medskip
\ifdefined\hypertarget\hypertarget{sourcepointer-4}{}\fi
\noindent\textit{Source reference.} Original lemma 089W: sources/stacks-chow/tags.txt, source lines 10511--10511. Navigation pointer for \textbackslash{}ref\{coherent-lemma-relative-affine-cohomology\} used inside the counted span (line 4629); the argument that consumes it is retained in the span.\par\medskip
\ifdefined\hypertarget\hypertarget{sourcepointer-5}{}\fi
\noindent\textit{Source reference.} Original proposition 02UZ: sources/stacks-chow/tags.txt, source lines 3750--3750. Navigation pointer for \textbackslash{}ref\{cohomology-proposition-vanishing-Noetherian\} used inside the counted span (line 4640); the argument that consumes it is retained in the span.\par\medskip
\ifdefined\hypertarget\hypertarget{sourcepointer-6}{}\fi
\noindent\textit{Source reference.} Original remark 0FVX: sources/stacks-chow/tags.txt, source lines 19827--19827. Navigation pointer for \textbackslash{}ref\{duality-remark-coherent-duality-proper-over-field\} used inside the counted span (line 4593); the argument that consumes it is retained in the span.\par\medskip
\ifdefined\hypertarget\hypertarget{sourcepointer-7}{}\fi
\noindent\textit{Source reference.} Original lemma 0FW2: sources/stacks-derham/pinned-original-source.tex, source lines 4368--4368. Navigation pointer for \textbackslash{}ref\{lemma-Garel-map-frobenius-smooth-char-p\} used inside the counted span (line 4616); the argument that consumes it is retained in the span.\par\medskip
\ifdefined\hypertarget\hypertarget{sourcepointer-8}{}\fi
\noindent\textit{Source reference.} Original lemma 0FW5: sources/stacks-derham/pinned-original-source.tex, source lines 4548--4548. Navigation pointer for \textbackslash{}ref\{lemma-bottom-part-degenerates\} used inside the counted span (line 4588); the argument that consumes it is retained in the span.\par\medskip
\ifdefined\hypertarget\hypertarget{sourcepointer-9}{}\fi
\noindent\textit{Source reference.} Original lemma 0FW4: sources/stacks-derham/pinned-original-source.tex, source lines 4496--4496. Navigation pointer for \textbackslash{}ref\{lemma-duality-hodge\} used inside the counted span (line 4588); the argument that consumes it is retained in the span.\par\medskip
\ifdefined\hypertarget\hypertarget{sourcepointer-10}{}\fi
\noindent\textit{Source reference.} Original lemma 0EY2: sources/stacks-chow/tags.txt, source lines 18677--18677. Navigation pointer for \textbackslash{}ref\{limits-lemma-descend-dimension-d\} used inside the counted span (line 4659); the argument that consumes it is retained in the span.\par\medskip
\ifdefined\hypertarget\hypertarget{sourcepointer-11}{}\fi
\noindent\textit{Source reference.} Original lemma 01ZM: sources/stacks-chow/tags.txt, source lines 2599--2599. Navigation pointer for \textbackslash{}ref\{limits-lemma-descend-finite-presentation\} used inside the counted span (line 4658); the argument that consumes it is retained in the span.\par\medskip
\ifdefined\hypertarget\hypertarget{sourcepointer-12}{}\fi
\noindent\textit{Source reference.} Original lemma 0C0C: sources/stacks-chow/tags.txt, source lines 15082--15082. Navigation pointer for \textbackslash{}ref\{limits-lemma-descend-smooth\} used inside the counted span (line 4659); the argument that consumes it is retained in the span.\par\medskip
\ifdefined\hypertarget\hypertarget{sourcepointer-13}{}\fi
\noindent\textit{Source reference.} Original lemma 081F: sources/stacks-chow/tags.txt, source lines 10215--10215. Navigation pointer for \textbackslash{}ref\{limits-lemma-eventually-proper\} used inside the counted span (line 4660); the argument that consumes it is retained in the span.\par\medskip
\ifdefined\hypertarget\hypertarget{sourcepointer-14}{}\fi
\noindent\textit{Source reference.} Original lemma 0E7D: sources/stacks-chow/tags.txt, source lines 17778--17778. Navigation pointer for \textbackslash{}ref\{limits-lemma-higher-direct-images-zero-above-dimension-fibre\} used inside the counted span (line 4682); the argument that consumes it is retained in the span.\par\medskip
\ifdefined\hypertarget\hypertarget{sourcepointer-15}{}\fi
\noindent\textit{Source reference.} Original lemma 01UX: sources/stacks-chow/tags.txt, source lines 2426--2426. Navigation pointer for \textbackslash{}ref\{morphisms-lemma-triangle-differentials\} used inside the counted span (line 4607); the argument that consumes it is retained in the span.\par\medskip
\ifdefined\hypertarget\hypertarget{sourcepointer-16}{}\fi
\noindent\textit{Source reference.} Original lemma 0BUG: sources/stacks-chow/tags.txt, source lines 14874--14874. Navigation pointer for \textbackslash{}ref\{varieties-lemma-proper-geometrically-reduced-global-sections\} used inside the counted span (line 4601); the argument that consumes it is retained in the span.\par\medskip
\ifdefined\hypertarget\hypertarget{sourcepointer-17}{}\fi
\noindent\textit{Source reference.} Original lemma 056T: sources/stacks-chow/tags.txt, source lines 6721--6721. Navigation pointer for \textbackslash{}ref\{varieties-lemma-smooth-geometrically-normal\} used inside the counted span (line 4600); the argument that consumes it is retained in the span.\par\medskip
\begin{lemma}
\label{lemma-top-part-degenerates}
Let $k$ be a field. Let $X$ be a smooth proper scheme over $k$
equidimensional of dimension $d$. The map
$$
\text{d} : H^d(X, \Omega^{d - 1}_{X/k}) \to H^d(X, \Omega^d_{X/k})
$$
is zero.
\end{lemma}

\begin{proof}
It is tempting to think this follows from a combination of
Lemmas \ref{lemma-bottom-part-degenerates} and \ref{lemma-duality-hodge}.
However this doesn't work because the maps $\mathcal{O}_X \to \Omega^1_{X/k}$
and $\Omega^{d - 1}_{X/k} \to \Omega^d_{X/k}$ are not $\mathcal{O}_X$-linear
and hence we cannot use the functoriality discussed in
Duality for Schemes, Remark
\ref{duality-remark-coherent-duality-proper-over-field}
to conclude the map in Lemma \ref{lemma-bottom-part-degenerates}
is dual to the one in this lemma.

\medskip\noindent
We may replace $X$ by a connected component of $X$. Hence we may assume
$X$ is irreducible. By
Varieties, Lemmas \ref{varieties-lemma-smooth-geometrically-normal} and
\ref{varieties-lemma-proper-geometrically-reduced-global-sections}
we see that $k' = H^0(X, \mathcal{O}_X)$ is a finite separable
extension $k'/k$. Since $\Omega_{k'/k} = 0$
(see for example Algebra, Lemma
\ref{algebra-lemma-characterize-separable-algebraic-field-extensions})
we see that $\Omega_{X/k} = \Omega_{X/k'}$
(see Morphisms, Lemma \ref{morphisms-lemma-triangle-differentials}).
Thus we may replace $k$ by $k'$ and assume that $H^0(X, \mathcal{O}_X) = k$.

\medskip\noindent
Assume $H^0(X, \mathcal{O}_X) = k$. We conclude that
$\dim H^d(X, \Omega^d_{X/k}) = 1$ by Lemma \ref{lemma-duality-hodge}.
Assume first that the characteristic of $k$ is a prime number $p$.
Denote $F_{X/k} : X \to X^{(p)}$ the relative Frobenius of $X$ over $k$;
please keep in mind the facts proved about this morphism in
Lemma \ref{lemma-Garel-map-frobenius-smooth-char-p}.
Consider the commutative diagram
$$
\xymatrix{
H^d(X, \Omega^{d - 1}_{X/k}) \ar[d] \ar[r] &
H^d(X^{(p)}, F_{X/k, *}\Omega^{d - 1}_{X/k}) \ar[d] \ar[r]_{\Theta^{d - 1}} &
H^d(X^{(p)}, \Omega^{d - 1}_{X^{(p)}/k}) \ar[d] \\
H^d(X, \Omega^d_{X/k}) \ar[r] &
H^d(X^{(p)}, F_{X/k, *}\Omega^d_{X/k}) \ar[r]^{\Theta^d} &
H^d(X^{(p)}, \Omega^d_{X^{(p)}/k})
}
$$
The left two horizontal arrows are isomorphisms as $F_{X/k}$ is finite, see
Cohomology of Schemes, Lemma \ref{coherent-lemma-relative-affine-cohomology}.
The right square commutes as $\Theta_{X^{(p)}/X}$ is a morphism of
complexes and $\Theta^{d - 1}$ is zero. Thus it suffices to show that
$\Theta^d$ is nonzero (because the dimension of the source of the map
$\Theta^d$ is $1$ by the discussion above). However, we know that
$$
\Theta^d : F_{X/k, *}\Omega^d_{X/k} \to \Omega^d_{X^{(p)}/k}
$$
is surjective and hence surjective after applying the right exact
functor $H^d(X^{(p)}, -)$ (right exactness by the vanishing of cohomology
beyond $d$ as follows from
Cohomology, Proposition \ref{cohomology-proposition-vanishing-Noetherian}).
Finally, $H^d(X^{(d)}, \Omega^d_{X^{(d)}/k})$ is nonzero for example because
it is dual to $H^0(X^{(d)}, \mathcal{O}_{X^{(p)}})$ by
Lemma \ref{lemma-duality-hodge} applied to $X^{(p)}$ over $k$.
This finishes the proof in this case.

\medskip\noindent
Finally, assume the characteristic of $k$ is $0$.
We can write $k$ as the filtered colimit of its finite type
$\mathbf{Z}$-subalgebras $R$. For one of these we can find a
cartesian diagram of schemes
$$
\xymatrix{
X \ar[d] \ar[r] & Y \ar[d] \\
\Spec(k) \ar[r] & \Spec(R)
}
$$
such that $Y \to \Spec(R)$ is smooth of relative dimension $d$ and proper.
See Limits, Lemmas \ref{limits-lemma-descend-finite-presentation},
\ref{limits-lemma-descend-smooth}, \ref{limits-lemma-descend-dimension-d}, and
\ref{limits-lemma-eventually-proper}.
The modules $M^{i, j} = H^j(Y, \Omega^i_{Y/R})$ are finite $R$-modules, see
Cohomology of Schemes, Lemma
\ref{coherent-lemma-proper-over-affine-cohomology-finite}.
Thus after replacing $R$ by a localization we may assume all of these
modules are finite free. We have
$M^{i, j} \otimes_R k = H^j(X, \Omega^i_{X/k})$
by flat base change (Cohomology of Schemes, Lemma
\ref{coherent-lemma-flat-base-change-cohomology}).
Thus it suffices to show that $M^{d - 1, d} \to M^{d, d}$
is zero. This is a map of finite free modules over a domain,
hence it suffices to find a dense set of primes $\mathfrak p \subset R$
such that after tensoring with $\kappa(\mathfrak p)$ we get zero.
Since $R$ is of finite type over $\mathbf{Z}$, we can take
the collection of primes $\mathfrak p$ whose residue field
has positive characteristic (details omitted). Observe that
$$
M^{d - 1, d} \otimes_R \kappa(\mathfrak p) =
H^d(Y_{\kappa(\mathfrak p)},
\Omega^{d - 1}_{Y_{\kappa(\mathfrak p)}/\kappa(\mathfrak p)})
$$
for example by Limits, Lemma
\ref{limits-lemma-higher-direct-images-zero-above-dimension-fibre}.
Similarly for $M^{d, d}$. Thus we see that
$M^{d - 1, d} \otimes_R \kappa(\mathfrak p) \to
M^{d, d} \otimes_R \kappa(\mathfrak p)$
is zero by the case of positive characteristic handled above.
\end{proof}

\end{document}
